\section{Test Methodology}
\label{sec:test-methodology}

To rigorously evaluate the architectural limitations of native Kubernetes scaling primitives against an Operator-driven approach, this thesis adopts a deterministic, A/B comparative methodology. By intentionally eschewing active traffic generators and dynamic metric-driven autoscalers, the experimental design isolates the fundamental data safety characteristics of stateful scale-in operations. This methodology distils the evaluation to the overarching contrast between Kubernetes' reverse-ordinal pod deletion and the topology-aware data migration orchestrated by a domain-specific operator (\cite{k8s:operator}; \cite{opstree:redisoperator}).

\subsection{Definition of Success}
A foundational aspect of this methodology is the establishment of a rigorous, state-aware metric for what constitutes a successful scaling event. In a stateful context, capacity reduction without data integration is functionally equivalent to system failure. Therefore, a successful scale-in is strictly defined by three absolute criteria: maintaining a healthy \texttt{cluster\_state:ok} status, ensuring that all \num{16384} hash slots remain actively covered by the surviving nodes, and demonstrating zero data loss as evidenced by the perfect retention of the dataset and the total absence of \texttt{CLUSTERDOWN} errors during post-transition read probes (\cite{redis:clusterspec}).

\subsection{A/B Comparative Structure}
The methodology relies on an A/B comparative analysis designed to eliminate confounding variables by enforcing identical physical infrastructure and initial cluster geometries (\cite{k8s:statefulset}). Both scenarios initialize a strictly configured four-node Redis cluster composed exclusively of master nodes, deliberately omitting replica nodes to ensure that any topological disruption immediately surfaces as measurable data loss rather than being temporarily absorbed by high-availability failover mechanisms.
\begin{itemize}
    \item \textbf{Scenario A (Native Primitives):} Simulates the mechanical action of a Horizontal Pod Autoscaler by directly decrementing the replica count of the underlying \texttt{StatefulSet}. This scenario tests the hypothesis that Kubernetes' native, topology-blind termination sequence inherently violates the data integrity requirements of a sharded database.
    \item \textbf{Scenario B (Operator Pattern):} Modifies the declarative \texttt{clusterSize} parameter of the \texttt{RedisCluster} Custom Resource. This scenario tests the hypothesis that encoding domain-specific orchestration knowledge into a Kubernetes controller enables the pre-emptive, topology-aware resharding necessary for a graceful cluster downscale.
\end{itemize}

\subsection{High Reproducibility and Empirical Objectivity}
Scientific reproducibility is enforced through strict execution hygiene. Every experimental iteration is preceded by a programmatic "scorched earth" initialization phase that systematically uninstalls all Helm releases, purges the designated namespaces, and deletes all persistent volume claims. This absolute state destruction guarantees that no residual configurations or orphaned data fragments from prior executions can contaminate the baseline measurements. 

Furthermore, by performing the scaling operations under strictly quiescent conditions---where a precisely enumerated dataset of \num{5000} keys is populated prior to scaling, and no concurrent external write load is applied---the methodology achieves absolute empirical objectivity. It prevents intermittent network contention, connection tracking saturation, and load-induced resharding interruptions from confounding the results, thereby yielding a pristine measurement of the underlying scaling mechanisms' inherent capabilities.

\vspace{2em}